\chapter{Конструкторский раздел}

В этом разделе описан модифицированный метод ведения диалога роботом Ф-2 с учётом персональных черт робота.

\section{Предлагаемый метод}

На рисунке~\ref{img:f2_speech_method} показана функциональная схема модифицированного метода ведения диалога роботом Ф-2 с учётом персональных черт робота. Модификация метода будет разработана на основе введения блока A3, представленного также
на рисунке~\ref{img:fpz} %Это НЕ ФПЗ, это блок А3
и имеющего 4 управляющих воздействия.
На рисунке~\ref{img:filter_method_specified} представлена детализация функциональной модели блока A3 предлагаемого метода в нотации IDEF0.

\includesvgimage
{f2_speech_method}
{f}
{H}
{1\textwidth}
{Модифицированный метод ведения диалога роботом Ф-2 с учётом персональных черт робота}

\includesvgimage
{fpz}
{f}
{H}
{1\textwidth}
{Функциональная модель модуля модификацииповедения на основе персональных черт}


\includesvgimage
{filter_method_specified}
{f}
{H}
{1\textwidth}
{Детализация модуля модификации поведения на основе персональных черт}

Базовый метод ведения диалога роботом Ф-2 в рассматриваемой части можно представить следующей последовательностью действий~\ref{F2arch}:
\begin{enumerate}
	\item получение предложения --- высказывания пользователя;
	\item построение его семантического представления;
	\item сравнение семантического представления с посылками сценариев;
	\item вычисление весов активированных сценариев;
	\item выбор сценария с наибольшим весом;
	\item выбор мультимодальной реакции из базы реакций;
	\item воспроизведение BML-кадра поведения.
\end{enumerate}

В модифицированном методе в данную последовательность между пунктами 4 и 5 предлагается добавить один блок --- модуль модификации поведения на основе персональных черт.
Входные данные модуля определяются следующим множеством:
\begin{equation}
	A = \{\langle d_i, w_i \rangle \mid i = 1 \dots N\},
\end{equation}
где

\begin{itemize}
	\item \(d_i\) --- активированный д-сценарий;
	\item \(w_i\) --- исходный вес активированного д-сценария \(d_i\);
	\item \(N\) --- число рассматриваемых д-сценариев.
\end{itemize}

На основе задаваемого профиля (характера) робота и правил переформирования набора сценариев модуль вычисляет множитель \(m_i\) для каждого активированного д-сценария.

Затем вес каждого активированного д-сценария пересчитывается по формуле:
\begin{equation}
	\label{eq:modified_weight}
	w'_i = clip(w_i \cdot m_i),
\end{equation}
где

\begin{itemize}
	\item \(w'_i\) --- модифицированный вес \(i\)-го активированного д-сценария;
	\item \(clip\) --- отсечение значения к допустимому диапазону весов сценариев.
\end{itemize}

Выходом модуля является модифицированный набор активированных д-сценариев:
\begin{equation}
	\label{eq:method_output}
	A' = \{\langle d_i, w'_i \rangle \mid i = 1 \dots N\},
\end{equation}
где

\begin{itemize}
	\item \(d_i\) --- активированный д-сценарий;
	\item \(w'_i\) --- вес активированного д-сценария \(d_i\) после применения модуля.
\end{itemize}

Данный набор передаётся в блок выбора сценария робота Ф-2.

Для данного модуля выделены 4 управляющих воздействия:
\begin{enumerate}
	\item характер, задаваемый персональными чертами;
	\item параметры отношения робота к собеседнику;
	\item параметры ситуации;
	\item правила переформирования набора сценариев.
\end{enumerate}

В основе метода лежит концепт жизненного пространства, предложенный К.~Левиным~\cite{grishina2018context}.
В статье Н.В.~Гришиной данный концепт рассматривается как единица описания жизненного контекста, объединяющая личность и ситуацию в целостную схему анализа поведения. При этом бытийный уровень контекста, связанный в~\cite{grishina2018context} с психологией человеческого бытия и экзистенциальной психологии, не входит в рамки данной работы. При этом т.~н. психология среды, изучающая когнитивные и поведенческие ответы людей
на характеристики среды (внешние воздействия на человека), отражена в аппарате когнитивного компонента робота Ф-2 (см.~блок А2 на рисунке~\ref{img:f2_speech_method}).

\section{Параметры предлагаемого метода}

Параметры предлагаемого метода включают следующие части:

\begin{itemize}
	\item профиль персональных черт робота;
	\item параметры отношения робота к собеседнику;
	\item параметры ситуации.
\end{itemize}

\subsection*{Профиль персональных черт робота}

Профиль персональных черт робота задаётся множеством термов проявления граней <<Большой пятёрки>>:
\begin{equation}
	\label{eq:personality_profile}
	P = \{\langle p_j, value_j \rangle \mid j = 1 \dots K\},
\end{equation}
где

\begin{itemize}
	\item \(p_j\) --- грань личности;
	\item \(value_j\) --- терм, описывающий способ проявления грани \(p_j\);
	\item \(K\) --- число используемых граней.
\end{itemize}

Термы используются потому, что одинаковая выраженность одной и той же черты у разных характеров может проявляться по-разному.

\subsection*{Параметры отношения робота к собеседнику}

Параметры отношения робота к собеседнику задаются множеством:
\begin{equation}
	\label{eq:interlocutor_relation}
	R = \{\langle r_l, value_l \rangle \mid l = 1 \dots L\},
\end{equation}
где

\begin{itemize}
	\item \(r_l\) --- параметр отношения робота к собеседнику;
	\item \(value_l\) --- значение параметра \(r_l\);
	\item \(L\) --- количество параметров отношения, используемых в конкретной конфигурации профиля.
\end{itemize}

Для каждого параметра отношения задаётся область значений:
\begin{equation}
	\label{eq:relation_domain}
	value_l \in D^R_l.
\end{equation}

Область \(D^R_l\) может быть задана термом или числом.
Если параметр отношения сразу задаётся термами, то \(D^R_l = T^R_l\), где \(T^R_l\) --- множество допустимых лингвистических значений параметра.
Если параметр задаётся числом, то \(D^R_l = X^R_l\), где \(X^R_l \subseteq \mathbb{R}\), а для интерпретации численного значения дополнительно задаётся терм-множество \(T^R_l\) и функции принадлежности:
\begin{equation}
	\label{eq:relation_membership}
	\mu^R_{l,t}: X^R_l \rightarrow [0; 1], \quad t \in T^R_l.
\end{equation}

Тем самым численный параметр отношения перед использованием в правилах переводится в степени принадлежности термам.
Заданный термом параметр может использоваться как чёткое значение или как распределение степеней принадлежности по нескольким термам.

\subsection*{Параметры ситуации}

Параметры ситуации описывают контекст беседы.
Они также задаются множеством:
\begin{equation}
	\label{eq:situation_context}
	S = \{\langle s_h, value_h \rangle \mid h = 1 \dots M\},
\end{equation}
где

\begin{itemize}
	\item \(s_h\) --- параметр ситуации;
	\item \(value_h\) --- значение параметра \(s_h\);
	\item \(M\) --- количество параметров ситуации, выбранных для конкретной ситуации.
\end{itemize}

Для каждого параметра ситуации задаётся область значений:
\begin{equation}
	\label{eq:situation_domain}
	value_h \in D^S_h.
\end{equation}

Область \(D^S_h\) может быть задана термом или числом.
Если параметр ситуации задаётся термами, то \(D^S_h = T^S_h\), где \(T^S_h\) --- множество допустимых лингвистических значений.
Если параметр задаётся числом, то \(D^S_h = X^S_h\), где \(X^S_h \subseteq \mathbb{R}\), а для него дополнительно задаются терм-множество \(T^S_h\) и функции принадлежности:
\begin{equation}
	\label{eq:situation_membership}
	\mu^S_{h,t}: X^S_h \rightarrow [0; 1], \quad t \in T^S_h.
\end{equation}

Таким образом, профиль \(P\), отношение к собеседнику \(R\) и ситуация \(S\) имеют одинаковую структурную форму: это множества параметров, значения которых далее используются в правилах перевзвешивания д-сценариев.

При беседе с конкретным собеседником отношение к нему принимается за постоянную величину, в перспективе могут применяться правила, изменяющие отношение к собеседнику.
Ситуация может меняться при беседе с одним собеседником, поэтому параметры ситуации вынесены в отдельную группу.
В рамках данной работы параметры ситуации и отношения к собеседнику считаются параметрами, задаваемыми внешней силой и не рассчитываемыми по входному предложению.

\section{Нечёткий вывод}

Нечёткий вывод выполняется над термами профиля персональных черт робота, параметрами отношения робота к собеседнику и параметрами ситуации.
Если параметр уже задан лингвистическим значением, он используется в правилах напрямую.
Если параметр задан числом, то перед применением правил он фаззифицируется с помощью функций принадлежности, заданных для соответствующего терм-множества.

Для каждого численного параметра \(x_z\) с областью значений \(X_z\) задаётся множество термов \(T_z\) и функции принадлежности:
\begin{equation}
	\label{eq:generic_membership}
	\mu_{z,t}: X_z \rightarrow [0; 1], \quad t \in T_z.
\end{equation}

В результате численное значение параметра заменяется набором степеней принадлежности термам:
\begin{equation}
	\label{eq:fuzzified_feature}
	fuzz(x_z) = \{\langle t, \mu_{z,t}(x_z) \rangle \mid t \in T_z\}.
\end{equation}

Правило модуля модификации поведения на основе персональных черт имеет вид:
\begin{equation}
	\label{eq:rule_template}
	IF\; P_r \land C_r \land R_r\; THEN\; d_i \rightarrow e_r,
\end{equation}
где

\begin{itemize}
	\item \(P_r\) --- условие по профилю персональных черт;
	\item \(C_r\) --- условие по параметру ситуации;
	\item \(R_r\) --- условие по параметру отношения робота к собеседнику;
	\item \(e_r\) --- терм влияния правила.
\end{itemize}

Эффект правила принимает одно из значений, приведённых в таблице~\ref{tbl:multipliers}.
В данной работе используется схема нечёткого вывода Сугено нулевого порядка~\cite{takagi1985fuzzy}.
В такой схеме следствие правила задаётся не выходной функцией принадлежности, а точечным числовым значением:
терм влияния правила \(e_r\) отображается в числовой множитель \(k_r\) по таблице~\ref{tbl:multipliers}: например, эффект <<усиление>> соответствует \(k_r = 1.25\), а эффект <<сильное усиление>> соответствует \(k_r = 1.50\).
Итоговый выход вычисляется как среднее значений \(k_r\), взвешенное силами срабатывания правил.

\begin{table}[H]
	\centering
	\caption{Множители изменения весов д-сценариев}
	\label{tbl:multipliers}
	\begin{tabularx}{\textwidth}{|X|c|}
		\hline
		\textbf{Терм влияния правила} & \textbf{Численное значение} \\
		\hline
		сильное ослабление & 0.50 \\
		\hline
		ослабление & 0.75 \\
		\hline
		без изменения & 1.00 \\
		\hline
		усиление & 1.25 \\
		\hline
		сильное усиление & 1.50 \\
		\hline
	\end{tabularx}
\end{table}

Сила срабатывания правила вычисляется как минимум степеней истинности его условий:
\begin{equation}
	\label{eq:rule_activation}
	\alpha_r = \min(\mu_1, \mu_2, \dots, \mu_q),
\end{equation}
где

\begin{itemize}
	\item \(\alpha_r\) --- сила срабатывания правила \(r\);
	\item \(\mu_j\) --- степень истинности \(j\)-го условия;
	\item \(q\) --- число условий в правиле.
\end{itemize}

Если условие допускает несколько термов одного параметра, степень истинности такого условия вычисляется как максимум степеней принадлежности допустимым термам.

Если на один д-сценарий влияет несколько правил, их результаты объединяются через взвешенное среднее по силам срабатывания правил:
\begin{equation}
	\label{eq:multiplier_aggregation}
	m_i =
	\frac{\sum\limits_{r \in R_i}\alpha_r k_r}
	{\sum\limits_{r \in R_i}\alpha_r},
\end{equation}
где

\begin{itemize}
	\item \(R_i\) --- множество правил, влияющих на \(i\)-й д-сценарий;
	\item \(k_r\) --- числовой множитель, соответствующий следствию правила \(r\).
\end{itemize}

Если для сценария не сработало ни одно правило, принимается \(m_i = 1.00\).

\clearpage

\section{Алгоритм работы блока вычисления коэффициентов перевзвешивания сценариев}


На рисунке~\ref{img:scenario_reweighting_algorithm} показана схема алгоритма работы блока вычисления коэффициентов перевзвешивания сценариев.

\includesvgimage
{scenario_reweighting_algorithm}
{f}
{H}
{0.8\textwidth}
{Схема алгоритма работы блока вычисления коэффициентов перевзвешивания сценариев}

\section{Особенности и ограничения разработанного метода}

Предлагаемый метод имеет следующие особенности:

\begin{itemize}
	\item в данном методе не выбирается реакция напрямую, а изменяются веса уже активированных д-сценариев;
	\item черты задаются термами проявления, способы проявления черты определяются параметрами ситуации и отношения робота к собеседнику.
\end{itemize}

Предлагаемый метод имеет следующие ограничения:

\begin{itemize}
	\item рассматриваются только д-сценарии;
	\item в качестве входного предложения для робота подается простое предложение (повествовательное или побудительное) на русском языке;
	\item робот ведет диалог только с одним собеседником.
\end{itemize}

\section{Проектирование ПО}

В рамках данной работы будет спроектирован и реализован лишь блок А3 ---
модуль модификации поведения на основе персональных черт.
Блоки А1, А2, А4
и А5 рассматриваются как внешние по отношению к разрабатываемому модулю.

Разрабатываемый модуль обращается к синтактико-семантическому модулю
робота, включающему блоки А1 и А2.
Данный модуль принимает на вход предложение
пользователя на естественном языке.

Диаграмма компонентов разрабатываемого блока А3 приведена на
рисунке~\ref{img:components_diagram}.

\includesvgimage
{components_diagram}
{f}
{H}
{1\textwidth}
{Диаграмма компонентов разрабатываемого блока А3}

Диаграмма последовательностей взаимодействия разрабатываемого блока А3 с
внешними компонентами робота Ф-2 приведена на рисунке~\ref{img:diagram}.

\includesvgimage
{diagram}
{f}
{H}
{1\textwidth}
{Диаграмма последовательностей взаимодействия компонентов блока А3 с внешними компонентами робота Ф-2}


\phantomsection
\section*{Выводы}
\addcontentsline{toc}{section}{Выводы}

Разработан модифицированный метод ведения диалога роботом Ф-2 с учётом персональных черт робота.
Определены параметры метода, вид правил, описан алгоритм работы блока вычисления коэффициентов перевзвешивания сценариев.
Определена структура ПО в виде диаграммы компонентов и взаимодействие компонентов в виде диаграммы последовательностей.
Изложены особенности и ограничения предлагаемого метода.
